% =========================================================================
% ALGORITMOS PRIMAL-DUAL DE CUSTO MÍNIMO
% =========================================================================
\chapter{Algoritmos Primal-Dual de Custo Mínimo}\label{sec:algoritmos_mcf}

\section{A Abordagem Primal: Cancelamento de Ciclos Negativos (Cycle Canceling)}\label{sec:cycle_canceling}

\paragraph{Princípio Teórico.}
O algoritmo \textit{Cycle Canceling} de Klein~\cite{klein1967primal} opera sobre a condição de otimalidade por ciclos negativos: se um fluxo é viável mas não é de custo mínimo, o digrafo residual contém ao menos um ciclo direcionado de custo estritamente negativo.

\paragraph{Dinâmica Operacional em Duas Fases.}
O método organiza-se em duas etapas:

\begin{enumerate}
    \item \textbf{Fase 1 --- Viabilidade Primal:} Computa-se um fluxo máximo viável de $s$ a $t$ na rede, desconsiderando os custos, via \lstinline{FlowNetwork} (Dinic ou Edmonds-Karp).
    
    \item \textbf{Fase 2 --- Otimização por Cancelamento de Ciclos:} Com o fluxo máximo inicial:
    \begin{itemize}
        \item Constrói-se o digrafo residual $D_f$ com capacidades $c_f$ e custos $w_f$.
        \item Detecta-se um ciclo direcionado negativo $C$ via Bellman-Ford ou SPFA ($\mathcal{O}(|V| \cdot |A|)$).
        \item Calcula-se a capacidade residual mínima $\delta = \min_{(u,v) \in C} c_f(u,v) > 0$.
        \item Aumenta-se o fluxo ao longo de $C$ por $\delta$ unidades via \lstinline{push_flow}.
    \end{itemize}
\end{enumerate}

Para isolar o ciclo negativo a partir de um vértice $v$ relaxado na $|V|$-ésima etapa do algoritmo de caminho mais curto, retrocedem-se $|V|$ passos no vetor de predecessores. Pelo Princípio da Casa dos Pombos, o nó obtido pertence garantidamente ao ciclo negativo, viabilizando a extração de $C$. A classe \lstinline{CycleCanceling} implementa esse procedimento (Código~\ref{lst:cycle_canceling}, Apêndice~\ref{ap:cycle_canceling}).

Cada iteração preserva a viabilidade do fluxo e o valor $|f|$, reduzindo o custo total em $\delta \cdot |w(C)|$. O algoritmo encerra quando $D_f$ não possui ciclos negativos. A Figura~\ref{fig:cycle_canceling_exec} ilustra a dinâmica.

\begin{figure}[!ht]
\centering
\begin{tikzpicture}[scale=1.1, thick]
\begin{scope}[xshift=0cm]
  \tikzsubcaption{(2.5,2.6)}{(a) Ciclo negativo em $D_f$ detectado}
  \node[source node] (s) at (0, 1) {$s$};
  \node[flow node] (u) at (2.5, 2) {$u$};
  \node[flow node] (v) at (2.5, 0) {$v$};
  \node[sink node] (t) at (5, 1) {$t$};

  \draw[flow edge] (s) -- node[edge lbl below, pos=0.55] {$f\!=\!4,\; w\!=\!4$} (u);
  \draw[flow edge] (u) -- node[edge lbl above] {$f\!=\!4,\; w\!=\!2$} (t);
  \draw[flow edge, opacity=0.3] (s) -- (v);
  \draw[flow edge, opacity=0.3] (v) -- (t);

  \draw[sat edge, line width=1.6pt] (s) -- node[edge lbl below] {$w\!=\!1$} (v);
  \draw[sat edge, line width=1.6pt] (v) -- node[edge lbl right] {$w\!=\!1$} (u);
  \draw[rev edge, line width=1.6pt, bend right=32] (u) to node[edge lbl above, text=red!80!black] {$w_f\!=\!{-4}$} (s);

  \node[font=\footnotesize, align=center, below=4pt] at (2.5,-0.6) {Ciclo $C = s \to v \to u \to s$ em $D_f$\\$w(C) = 1 + 1 + (-4) = -2 < 0$\\Gargalo $\delta=4 \implies$ cancelamento reduz custo};
\end{scope}

\begin{scope}[xshift=7.5cm]
  \tikzsubcaption{(2.5,2.6)}{(b) Após cancelamento de $\delta=4$ unidades}
  \node[source node] (s2) at (0, 1) {$s$};
  \node[flow node] (u2) at (2.5, 2) {$u$};
  \node[flow node] (v2) at (2.5, 0) {$v$};
  \node[sink node] (t2) at (5, 1) {$t$};

  \draw[flow edge, opacity=0.3] (s2) -- (u2);
  \draw[flow edge] (u2) -- node[edge lbl above] {$f\!=\!4,\; w\!=\!2$} (t2);
  \draw[tree edge] (s2) -- node[edge lbl below] {$f\!=\!4,\; w\!=\!1$} (v2);
  \draw[tree edge] (v2) -- node[edge lbl right] {$f\!=\!4,\; w\!=\!1$} (u2);

  \node[font=\footnotesize, align=center, below=4pt] at (2.5,-0.6) {Fluxo líquido $|f|=4$ permanece inalterado\\Custo total reduz de $24$ para $16$};
\end{scope}
\end{tikzpicture}
\caption{Dinâmica do algoritmo Cycle-Canceling: (a) detecção do ciclo negativo em $D_f$. (b) Envio de $\delta=4$ unidades, reduzindo o custo sem alterar o fluxo líquido.}
\label{fig:cycle_canceling_exec}
\end{figure}

\paragraph{Complexidade Assintótica.}
Para capacidades inteiras limitadas por $C$ e custos por $W$, o custo total é no máximo $|A| \cdot C \cdot W$. Como cada cancelamento reduz o custo em pelo menos 1 unidade inteira, ocorrem no máximo $\mathcal{O}(|A| \cdot C \cdot W)$ iterações. Com custo $\mathcal{O}(|V| \cdot |A|)$ por iteração de Bellman-Ford, a complexidade total no pior caso é $\mathcal{O}(|V| \cdot |A|^2 \cdot C \cdot W)$~\cite{ahuja1993network}.


\section{A Abordagem Dual: Caminhos de Custo Mínimo (Successive Shortest Path)}\label{sec:ssp}

O algoritmo \textit{Successive Shortest Path} (SSP)~\cite{ahuja1993network, schrijver2003combinatorial} adota a estratégia dual: mantém a ausência de ciclos negativos em $D_f$ como invariante contínua e aumenta o fluxo líquido por caminhos de menor custo até atingir a viabilidade.

\subsection{Invariante Dual e Aumentos Sucessivos}

Partindo do fluxo nulo ($f \equiv 0$), o digrafo residual inicial $D_0$ não possui ciclos negativos (assumindo custos iniciais não negativos em ciclos).

A cada iteração:
\begin{enumerate}
    \item Determina-se o caminho simples $P$ de menor custo residual de $s$ a $t$ em $D_f$.
    \item Calcula-se o gargalo $\delta = \min_{(u,v) \in P} c_f(u,v) > 0$.
    \item Aumenta-se o fluxo em $\delta$ ao longo de $P$ via \lstinline{push_flow}.
\end{enumerate}

\paragraph{Preservação da Invariante Dual.}
O aumento ao longo do caminho de custo mínimo não introduz ciclos de custo negativo em $D_f$. Se um novo ciclo $C$ contivesse arcos reversos $(v,u)$ com $(u,v) \in P$, seu custo residual seria $w_{f'}(v,u) = -w_f(u,v) = d(u) - d(v)$. Como para os demais arcos vale $w_f(x,y) \ge d(y) - d(x)$ pela desigualdade triangular, o custo total do ciclo satisfaz $w_{f'}(C) \ge \sum_C [d(y) - d(x)] = 0$, assegurando a ausência de ciclos negativos.


\subsection{Variante com SPFA}\label{sec:ssp_spfa}

Como arcos reversos adquirem custos negativos $w_f(v,u) = -w(u,v) < 0$, o algoritmo de Dijkstra não pode ser aplicado diretamente sobre $w_f$.

A classe \lstinline{SuccessiveShortest} (Código~\ref{lst:successive_shortest}, Apêndice~\ref{ap:successive_shortest}) utiliza o \textit{Shortest Path Faster Algorithm} (SPFA), variante de Bellman-Ford com fila de vértices relaxados.

Com custo $\mathcal{O}(|V| \cdot |A|)$ por busca e até $F = |f^*|$ aumentos unitários, a complexidade total é:
\begin{equation}
    \mathcal{O}(|V| \cdot |A| \cdot F)
\end{equation}

\begin{figure}[!ht]
\centering
\begin{tikzpicture}[scale=1.05, thick]
\begin{scope}[xshift=0cm]
  \tikzsubcaption{(2.5,2.8)}{(a) Grafo residual com arco reverso negativo}
  \node[source node] (s) at (0, 1.0) {$s$};
  \node[flow node] (u) at (2.5, 2.0) {$u$};
  \node[flow node] (v) at (2.5, 0.0) {$v$};
  \node[sink node] (t) at (5.0, 1.0) {$t$};

  \draw[flow edge] (s) -- node[edge lbl below, pos=0.45] {$w_f\!=\!1$} (v);
  \draw[flow edge] (v) -- node[edge lbl below] {$w_f\!=\!2$} (t);
  \draw[flow edge] (u) -- node[edge lbl right] {$w_f\!=\!3$} (t);
  \draw[rev edge, line width=1.4pt] (u) -- node[edge lbl left, text=red!80!black] {$w_f\!=\!{-4}$} (v);
  \draw[flow edge, dashed, gray] (s) -- (u);

  \node[font=\footnotesize, align=center, below=5pt] at (2.5,-0.7) {Arco reverso $u \to v$ com $w_f = -4 < 0$};
\end{scope}

\begin{scope}[xshift=7.5cm]
  \tikzsubcaption{(2.5,2.8)}{(b) Dinâmica da fila SPFA}
  \node[source node] (s2) at (0, 1.0) {$s$};
  \node[active node] (u2) at (2.5, 2.0) {$u$};
  \node[active node] (v2) at (2.5, 0.0) {$v$};
  \node[sink node] (t2) at (5.0, 1.0) {$t$};

  \draw[tree edge] (s2) -- node[edge lbl above, pos=0.45, text=blue!80!black] {$d(u)\!=\!2$} (u2);
  \draw[flow edge] (s2) -- node[edge lbl below, pos=0.45] {$d(v)\!=\!1$} (v2);
  \draw[sat edge, line width=1.5pt] (u2) -- node[edge lbl left, text=red!80!black] {$w_f\!=\!{-4}$} (v2);
  \draw[flow edge] (v2) -- node[edge lbl below] {$d(t)\!=\!1$} (t2);

  \node[font=\footnotesize, align=center, below=5pt] at (2.5,-0.7) {Relaxação: $d(v)$ reduz de $1$ para $2 - 4 = -2$\\Vértice $v$ é reinserido na fila};
\end{scope}
\end{tikzpicture}
\caption{Mecânica do SPFA em caminhos aumentantes de custo mínimo: relaxação de arcos com custo residual negativo.}
\label{fig:spfa_mcf}
\end{figure}


\subsection{Variante com Potenciais Nodais e Dijkstra}\label{sec:ssp_dijkstra}

A técnica de potenciais nodais~\cite{edmonds1972theoretical} transforma os custos em valores não negativos $\bar{w}_\pi(u,v) = w(u,v) - \pi(u) + \pi(v) \ge 0$, viabilizando o uso de Dijkstra.

A cada iteração, calculam-se as distâncias mínimas $d(v)$ a partir de $s$ sob $\bar{w}_\pi$. Atualizando $\pi'(v) = \pi(v) - d(v)$, o novo custo reduzido satisfaz:
\begin{equation}
    \bar{w}_{\pi'}(u,v) = \bar{w}_\pi(u,v) + d(u) - d(v) \ge 0
\end{equation}
Para arcos reversos criados ao longo do caminho de menor custo, vale $d(v) = d(u) + \bar{w}_\pi(u,v)$, de modo que $\bar{w}_{\pi'}(v,u) = -\bar{w}_{\pi'}(u,v) = 0 \ge 0$.

A classe \lstinline{SuccessiveShortestDijkstra} (Código~\ref{lst:successive_shortest_dijkstra}, Apêndice~\ref{ap:successive_shortest_dijkstra}) implementa essa estratégia com complexidade:
\begin{equation}
    \mathcal{O}(F \cdot |A| \log |V|)
\end{equation}

\begin{figure}[!ht]
\centering
\begin{tikzpicture}[thick]
    \node[source node] (s) at (0, 1.5) {$s$};
    \node[flow node] (u) at (3.5, 3.0) {$u$};
    \node[flow node] (v) at (3.5, 0.0) {$v$};
    \node[sink node] (t) at (7.0, 1.5) {$t$};

    \node[font=\footnotesize, blue!80!black, left=3pt] at (s.west) {$\pi(s)\!=\!0$};
    \node[font=\footnotesize, blue!80!black, above=3pt] at (u.north) {$\pi(u)\!=\!2$};
    \node[font=\footnotesize, blue!80!black, below=3pt] at (v.south) {$\pi(v)\!=\!1$};
    \node[font=\footnotesize, blue!80!black, right=3pt] at (t.east) {$\pi(t)\!=\!5$};

    \draw[tree edge] (s) -- node[edge lbl below, text=blue!80!black, pos=0.55] {$c_f\!=\!3, \;\bar{w}\!=\!0$} (u);
    \draw[tree edge] (u) -- node[edge lbl above, text=blue!80!black, pos=0.5] {$c_f\!=\!2, \;\bar{w}\!=\!0$} (t);

    \draw[flow edge] (s) -- node[edge lbl below, pos=0.5] {$c_f\!=\!4, \;\bar{w}\!=\!1$} (v);
    \draw[flow edge] (v) -- node[edge lbl below, pos=0.5] {$c_f\!=\!1, \;\bar{w}\!=\!2$} (t);
    \draw[flow edge] (u) -- node[edge lbl, right=3pt, pos=0.5] {$c_f\!=\!2, \;\bar{w}\!=\!1$} (v);

    \draw[rev edge, bend right=38] (u) to node[edge lbl above, text=red!80!black, pos=0.5] {reverso: $\bar{w}'\!=\!0$} (s);
\end{tikzpicture}
\caption{Iteração de \textit{Successive Shortest Path} com potenciais nodais $\pi$. O caminho mais curto possui $\bar{w}_\pi = 0$, e arcos reversos gerados preservam $\bar{w}' = 0 \ge 0$.}
\label{fig:ssp_iteracao}
\end{figure}


\subsection{Equivalência com o Algoritmo Húngaro (Kuhn-Munkres)}

O Algoritmo Húngaro~\cite{kuhn1955hungarian, munkres1957algorithms} para o problema de atribuição bipartida ponderada equivale à execução do método \textit{Successive Shortest Path} com potenciais nodais sobre redes unitárias com balanceamento $b(v) \in \{+1, -1\}$. Os ajustes de potenciais da matriz de custos correspondem à manutenção da não negatividade dos custos reduzidos $\bar{w}_\pi \ge 0$.


\subsection{Comparação entre as Variantes}

\begin{itemize}
    \item \textbf{SPFA:} Opera diretamente sobre custos residuais sem manter o vetor de potenciais $\pi$, com menor complexidade de código, indicada para redes esparsas com poucos aumentos.
    \item \textbf{Dijkstra com Potenciais:} Requer a manutenção dos potenciais nodais $\pi$, mas assegura tempo de pior caso $\mathcal{O}(F \cdot |A| \log |V|)$ com fila de prioridades, apresentando menor tempo de CPU em redes densas.
\end{itemize}

